\chapter*{Uso de los talleres}\addcontentsline{toc}{chapter}{Uso de los talleres}
Cada taller incluye objetivo, código, resultado de contraste y trabajo propuesto. Los archivos editables están en \texttt{workshops/}; se ejecutan con Wolfram Language. El material de \texttt{companion/} pertenece a las prácticas anteriores del manual. Los scripts terminan con error si falla su comprobación. El informe del repositorio registra su ejecución real, versión y salida.

Todos los cálculos usan los datos hipotéticos del banco. Un cálculo simbólico o numérico que reproduce una ecuación no demuestra por sí solo que esa ecuación describa un experimento.


\chapter{Unidades antes de calcular}

\objetivos{Convierte caudal y radio a SI antes de evaluar la caída de presión. Repite el cálculo conservando cantidades con unidades para detectar una confusión entre milímetros y metros.}

\textbf{Modelo de referencia.} MF-P17.\par

\section{Procedimiento y código}

Antes de ejecutar, escribe las unidades esperadas y el control que decidirá si el cálculo coincide con el modelo.

\lstinputlisting{../workshops/w01.wl}

\section{Resultado de contraste}

$Q=10^{-6}\,\mathrm{m^3/s}$ y $\Delta p\simeq24446.2\,\mathrm{Pa}$. Introducir el radio en milímetros como si fueran metros modifica el resultado por un factor $10^{12}$.

\section{Trabajo propuesto}

Cambia el radio conservando el caudal. Explica la potencia cuarta. Introduce deliberadamente una unidad incompatible y observa el diagnóstico; no lo conviertas en un valor físico.

\studiumchaptercheck

\chapter{Integración temporal de una resistencia cuadrática}

\objetivos{Compara un integrador de ecuaciones diferenciales con la solución separada. La coincidencia controla implementación y tolerancias, sin validar la ley de fuerza frente a medidas.}

\textbf{Modelo de referencia.} MF-P02.\par

\section{Procedimiento y código}

Antes de ejecutar, escribe las unidades esperadas y el control que decidirá si el cálculo coincide con el modelo.

\lstinputlisting{../workshops/w02.wl}

\section{Resultado de contraste}

$v(6)=3{,}076923\,\mathrm{m/s}$. El error absoluto en 101 puntos debe quedar por debajo de $10^{-7}\,\mathrm{m/s}$ con la tolerancia propuesta.

\section{Trabajo propuesto}

Repite con tolerancias más laxas. Compara error de integración con error del modelo. Comprueba el límite de resistencia nula de la solución analítica.

\studiumchaptercheck

\chapter{Operadores vectoriales y aceleración}

\objetivos{Calcula divergencia, rotacional y aceleración convectiva del campo de remanso sin confundir un campo estacionario con partículas sin aceleración.}

\textbf{Modelo de referencia.} MF-P10.\par

\section{Procedimiento y código}

Antes de ejecutar, escribe las unidades esperadas y el control que decidirá si el cálculo coincide con el modelo.

\lstinputlisting{../workshops/w03.wl}

\section{Resultado de contraste}

$\nabla\cdot\boldsymbol u=0$, $\boldsymbol\omega=0$ y $\boldsymbol a=(0{,}108,0{,}072,0)\,\mathrm{m/s^2}$ en el punto dado.

\section{Trabajo propuesto}

Sustituye el campo por una cizalla simple. Explica por qué puede aparecer rotacional sin aceleración. Reconstruye la matriz simétrica y antisimétrica del gradiente.

\studiumchaptercheck

\chapter{Depósito: continuidad y evolución de nivel}

\objetivos{Integra el nivel hasta el tiempo analítico necesario para alcanzar el nivel final. Evita prolongar el modelo a niveles negativos.}

\textbf{Modelo de referencia.} MF-P34.\par

\section{Procedimiento y código}

Antes de ejecutar, escribe las unidades esperadas y el control que decidirá si el cálculo coincide con el modelo.

\lstinputlisting{../workshops/w04.wl}

\section{Resultado de contraste}

El nivel final numérico debe coincidir con $0{,}5\,\mathrm m$ dentro de $10^{-7}\,\mathrm m$. El caudal inicial no permanece constante durante el vaciado.

\section{Trabajo propuesto}

Duplica la sección del depósito. Repite con otro coeficiente dado y explica el cambio temporal. No interpretes ajustar el coeficiente como una validación sin datos.

\studiumchaptercheck

\chapter{Discretización espacial de Poiseuille plano}

\objetivos{Resuelve la ecuación de contorno mediante una matriz tridiagonal. En este caso el perfil exacto es cuadrático: la segunda diferencia centrada lo reproduce exactamente salvo redondeo.}

\textbf{Modelo de referencia.} MF-P16 con pared superior fija.\par

\section{Procedimiento y código}

Antes de ejecutar, escribe las unidades esperadas y el control que decidirá si el cálculo coincide con el modelo.

\lstinputlisting{../workshops/w05.wl}

\section{Resultado de contraste}

El error debe ser próximo al redondeo. Este resultado no mide orden de convergencia: la derivada cuarta de una parábola es cero. El taller 7, de difusión, sí permite observar error espacial no nulo.

\section{Trabajo propuesto}

Comprueba ambos valores de pared y el máximo central. Añade la contribución lineal de una pared móvil. Distingue residuo discreto de error respecto de la solución continua.

\studiumchaptercheck

\chapter{Ajuste de una ley de resistencia sintética}

\objetivos{Usa mínimos cuadrados con cuatro puntos construidos de una misma ley. La recuperación exacta de sus parámetros es un control algebraico, no evidencia experimental.}

\textbf{Modelo de referencia.} MF-P20.\par

\section{Procedimiento y código}

Antes de ejecutar, escribe las unidades esperadas y el control que decidirá si el cálculo coincide con el modelo.

\lstinputlisting{../workshops/w06.wl}

\section{Resultado de contraste}

$a=0{,}09\,\mathrm{kg/s}$ y $b=0{,}03\,\mathrm{kg/m}$. El residuo casi nulo se debe a que los datos sintéticos obedecen exactamente a la ley propuesta.

\section{Trabajo propuesto}

Añade una perturbación conocida a una fuerza y compara parámetros y residuos. Reserva un punto para prueba. Sin un modelo de incertidumbre no atribuyas significado estadístico a los decimales.

\studiumchaptercheck

\chapter{Difusión: convergencia de la tasa modal}

\objetivos{Compara la tasa de decaimiento continua del seno con la tasa de una segunda diferencia centrada. Duplica sucesivamente el número de intervalos y estima el orden de error espacial.}

\textbf{Modelo de referencia.} MF-P24.\par

\section{Procedimiento y código}

Antes de ejecutar, escribe las unidades esperadas y el control que decidirá si el cálculo coincide con el modelo.

\lstinputlisting{../workshops/w07.wl}

\section{Resultado de contraste}

Las razones de errores se aproximan a cuatro, correspondientes a segundo orden espacial. Esta comparación de tasas no introduce un integrador temporal ni mide su error.

\section{Trabajo propuesto}

Representa el error frente a la holgura de malla en escala logarítmica. Deduce el orden con una expansión de seno. Diferencia convergencia espacial de estabilidad y convergencia temporal.

\studiumchaptercheck
